\section{符号说明}

\begin{table}[H]
  \centering
  \caption{问题一符号说明}
  \label{tab:symbols}
  \begin{tabularx}{0.9\textwidth}{@{}cXc@{}}
    \toprule
    符号 & 含义 & 单位 \\
    \midrule
    $n$ & 检测点的数量 & 个 \\
    $S_i=(x_i,y_i)$ & 第 $i$ 个检测点的二维位置坐标 & $\mathrm{m}$ \\
    $\theta_i$ & 第 $i$ 个检测点在位置 $S_i$ 处测得的示向度 & $^\circ$ \\
    $\delta$ & 示向度的最大误差，本文取 $1^\circ$ & $^\circ$ \\
    $C_i$ & 第 $i$ 次观测所确定的角度误差区域 & -- \\
    $\Omega$ & 所有角度误差区域的公共部分，即定位区域 & -- \\
    $O_0$ & 半径为 $1800\ \mathrm{m}$ 的目标圆域圆心 & m \\
    $\overline B(A,\rho)$ & 以 $A$ 为圆心、$\rho$ 为半径的闭圆域 & -- \\
    $K$ & 同时考虑角度、接收距离上界和目标圆域的可行域 & -- \\
    $\widehat K$ & 问题二中以外接多边形代替圆域后得到的外近似 & -- \\
    $m$ & 定位凸多边形的顶点数 & 个 \\
    $V_j$ & 定位凸多边形的第 $j$ 个顶点 & m \\
    $d(P,Q)$ & 平面内两点 $P,Q$ 之间的欧氏距离 & m \\
    $D$ & 定位区域的直径 & m \\
    $V_p,V_q$ & 一组取得直径的端点 & m \\
    $O,r$ & 直径圆的圆心和半径 & m \\
    \bottomrule
  \end{tabularx}
\end{table}

\begin{table}[H]
  \centering
  \caption{问题二符号说明}
  \label{tab:q2-symbols}
  \begin{tabularx}{0.94\textwidth}{@{}cXc@{}}
    \toprule
    符号 & 含义 & 单位 \\
    \midrule
    $K_1$ & 第一次检测后同时考虑示向误差、目标圆域和接收距离上界的可行域 & -- \\
    $\widehat K_1$ & 以固定边数外切正多边形代替圆域后得到的 $K_1$ 保守外近似 & -- \\
    $X=(x,y)$ & 干扰源的未知位置坐标 & $\mathrm{m}$ \\
    $S_2=s$ & 待选择的第二检测点位置坐标 & $\mathrm{m}$ \\
    $R_{\min},R_{\max}$ & 全向干扰源有效接收半径的下界与上界 & $\mathrm{m}$ \\
    $N$ & 圆域固定外切正多边形的边数，本文取 $24$ & 边 \\
    $\widehat B_{\mathrm{out}}^{(N)}(A,R)$ & 圆盘 $\overline B(A,R)$ 的固定 $N$ 边外切正多边形 & -- \\
    $\Delta_{\mathrm{out}}(R,N)$ & 固定外切正多边形相对原圆盘的最大径向外扩量 & $\mathrm{m}$ \\
    $\mathcal F_g$ & 第二检测点的保证接收域 & -- \\
    $\mathcal F_p$ & 第二检测点的可能接收域 & -- \\
    $\rho(A)$ & 集合 $A$ 的定位不确定性半径，即最小包围圆半径 & $\mathrm{m}$ \\
    $\mathcal S_A,\mathcal E$ & 有限源位置场景集与有限示向误差场景集 & -- \\
    $R_{\mathrm{wc}}(s)$ & 二次定位不确定性半径的有限场景最坏预测值 & $\mathrm{m}$ \\
    $T(s)$ & 移动至 $s$ 并完成检测的行动时间 & $\mathrm{s}$ \\
    $\lambda_R$ & 综合评分中平衡时间与定位不确定性的权重超参数 & $\mathrm{s/m}$ \\
    $J(s)$ & 有限离散候选点 $s$ 的综合评分值 & $\mathrm{s}$ \\
    $\mathcal C_A,\mathcal C_A^g$ & 方案一的有限离散候选集及其中的保证接收子集 & -- \\
    $\Phi_{\mathrm{rob}}(s)$ & 对有限源位置场景取最小值的单位时间 Fisher 信息指标 & $\mathrm{m}^{-4}\mathrm{s}^{-1}$ \\
    $s_A,s_F,s^*$ & 离散方案代表点、连续方案代表点及最终推荐点 & $\mathrm{m}$ \\
    $\mathcal N_{TR}$ & 按行动时间与定位不确定性半径确定的非支配候选集 & -- \\
    \bottomrule
  \end{tabularx}
\end{table}

